\par\Needspace{6\baselineskip}\medskip
\hypertarget{native-102-article077}{}\pdfbookmark[2]{77.}{article-077}
\providecommand{\DailyMotionSign}{\mathord{\text{\itshape 7}}}
\begin{center}77.\end{center}
Iunctis iis, quae in art. praec. evoluta sunt, cum iis quae in artt. 15, 16, 20, 27, 28 pro singulis sectionum conicarum generibus tradidimus, omnia praesto erunt, quae ad calculum variationum differentialium loco geocentrico a variationibus singulorum elementorum inductarum requiruntur. Ad maiorem illustrationem horum praeceptorum exemplum supra in artt. 13, 14, 51, 63, 65 tractatum resumemus. Ac primo quidem ad normam art. praec. $\mathrm{d}l$ et $\mathrm{d}b$ per $\mathrm{d}r$, $\mathrm{d}u$, $\mathrm{d}i$, $\mathrm{d}\AscendingNode$ exprimemus, qui calculus ita se habet:
\begin{center}
\begin{tabular}{@{}ll@{}}
\makebox[17em][l]{$\log\operatorname{tang}\omega$\hfill .\,.\,.\,.} & $8{,}40113$\\[.4em]
\makebox[17em][l]{$\log\cos i$\hfill .\,.\,.\,.\,.} & $9{,}98853$\\[.4em]
\hline
\makebox[17em][l]{$\log\operatorname{tang}M$\hfill .\,.\,.} & $8{,}41260$\\[.4em]
$M=$ & $1^0\,28'\,52''$\\[.4em]
$M-u=$ & $165\quad17\quad8$\\[.4em]
\end{tabular}
\end{center}
\begin{center}
\begin{tabular}{@{}ll@{}}
\makebox[17em][l]{$\log\sin\omega$\hfill .\,.\,.\,.\,.} & $8{,}40099_n$\\[.4em]
\makebox[17em][l]{$\log\operatorname{tang}i$\hfill .\,.\,.\,.} & $9{,}36723$\\[.4em]
\hline
\makebox[17em][l]{$\log\operatorname{tang}N$\hfill .\,.\,.} & $7{,}76822_n$\\[.4em]
$N=$ & $179^0\,39'\,50''$\\[.4em]
$N-b=$ & $186\quad1\quad45$\\[.4em]
\end{tabular}
\end{center}
\begin{center}
\begin{tabular}{@{}ll@{}}
\makebox[17em][l]{$\log\operatorname{tang}(M-u)$\hfill .\,.\,.} & $9{,}41932_n$\\[.4em]
\makebox[17em][l]{$\log\cos\omega\sin i$\hfill .\,.\,.\,.\,.} & $9{,}35562_n$\\[.4em]
\hline
\makebox[17em][l]{$\log\operatorname{tang}P$\hfill .\,.\,.\,.\,.\,.\,.\,.} & $0{,}06370$\\[.4em]
$P=$ & $49^0\,11'\,13''$\\[.4em]
$N-b-P=$ & $136\quad50\quad32$\\[.4em]
\end{tabular}
\end{center}
\SourceLeaf{103}
\begin{center}
\begin{tabular}{@{}ll@{}}
\multicolumn{2}{c}{I$^{*}$}\\[.4em]
\makebox[17em][l]{$\log\sin(L-l)$\hfill .\,.\,.} & $9{,}72125$\\[.4em]
\makebox[17em][l]{$\log R$\hfill .\,.\,.\,.\,.\,.\,.\,.\,.\,.} & $9{,}99810$\\[.4em]
\makebox[17em][l]{$\mathrm{C.}\log\Delta'$\hfill .\,.\,.\,.\,.\,.\,.} & $9{,}92027$\\[.4em]
\hline
\makebox[17em][l]{$({}^{*})$\hfill .\,.\,.\,.\,.\,.\,.\,.\,.\,.\,.\,.} & $9{,}63962$\\[.4em]
\makebox[17em][l]{$\mathrm{C.}\log r$\hfill .\,.\,.\,.\,.\,.\,.\,.} & $9{,}67401$\\[.4em]
\hline
\makebox[17em][l]{$\log\left(\frac{\mathrm{d}l}{\mathrm{d}r}\right)$\hfill .\,.\,.\,.\,.\,.\,.} & $9{,}31363$\\[.4em]
\end{tabular}
\end{center}
\begin{center}
\begin{tabular}{@{}ll@{}}
\multicolumn{2}{c}{II$^{**}$}\\[.4em]
\makebox[17em][l]{$({}^{*})$\hfill .\,.\,.\,.\,.\,.\,.\,.\,.\,.\,.} & $9{,}63962$\\[.4em]
\makebox[17em][l]{$\log\operatorname{cot}(M-u)$\hfill .} & $0{,}58068_n$\\[.4em]
\hline
\makebox[17em][l]{$\log\left(\frac{\mathrm{d}l}{\mathrm{d}u}\right)$\hfill .\,.\,.\,.\,.\,.} & $0{,}22030$\\[.4em]
\end{tabular}
\end{center}
\begin{center}
\begin{tabular}{@{}ll@{}}
\multicolumn{2}{c}{III$^{*}$}\\[.4em]
\makebox[17em][l]{$\log\cos\omega$\hfill .\,.\,.\,.\,.\,.\,.\,.\,.\,.\,.} & $9{,}99986_n$\\[.4em]
\makebox[17em][l]{$\log\operatorname{tang}b$\hfill .\,.\,.\,.\,.\,.\,.\,.\,.\,.} & $9{,}04749_n$\\[.4em]
\hline
\makebox[17em][l]{$\log\left(\frac{\mathrm{d}l}{\mathrm{d}i}\right)$\hfill .\,.\,.\,.\,.\,.\,.\,.\,.\,.\,.} & $9{,}04735_n$\\[.4em]
\end{tabular}
\end{center}
\begin{center}
\begin{tabular}{@{}ll@{}}
\multicolumn{2}{c}{IV}\\[.4em]
\makebox[17em][l]{$\log\frac{R}{\Delta'}$\hfill .\,.\,.\,.\,.\,.\,.\,.\,.\,.} & $9{,}91837$\\[.4em]
\makebox[17em][l]{$\log\cos(L-l)$\hfill .\,.} & $9{,}92956$\\[.4em]
\hline
\makebox[17em][l]{$({}^{**})$\hfill .\,.\,.\,.\,.\,.\,.\,.\,.\,.\,.} & $9{,}84793$\\[.4em]
\multicolumn{2}{@{}l@{}}{$=\log\left(\frac{\mathrm{d}l}{\mathrm{d}\AscendingNode}-1\right)$}\\[.4em]
\end{tabular}
\end{center}
\begin{center}
\begin{tabular}{@{}ll@{}}
\multicolumn{2}{c}{V$^{*}$}\\[.4em]
\makebox[17em][l]{$({}^{**})$\hfill .\,.\,.\,.\,.\,.\,.\,.\,.\,.} & $9{,}84793$\\[.4em]
\makebox[17em][l]{$\log\sin b\cos b$\hfill .\,.} & $9{,}04212_n$\\[.4em]
\makebox[17em][l]{$\mathrm{C.}\log r$\hfill .\,.\,.\,.\,.\,.\,.} & $9{,}67401$\\[.4em]
\hline
\makebox[17em][l]{$\log\left(\frac{\mathrm{d}b}{\mathrm{d}r}\right)$\hfill .\,.\,.\,.\,.\,.} & $8{,}56406$\\[.4em]
\end{tabular}
\end{center}
\begin{center}
\begin{tabular}{@{}ll@{}}
\multicolumn{2}{c}{VI$^{**}$}\\[.4em]
\makebox[17em][l]{$\log\frac{r}{\Delta}$\hfill .\,.\,.\,.\,.\,.\,.\,.\,.\,.\,.\,.\,.} & $0{,}24357$\\[.4em]
\makebox[17em][l]{$\log\sin(M-u)$\hfill .\,.\,.\,.\,.} & $9{,}40484$\\[.4em]
\makebox[17em][l]{$\log\cos(N-b-P)$\hfill .\,.} & $9{,}86301_n$\\[.4em]
\makebox[17em][l]{$\mathrm{C.}\log\sin P$\hfill .\,.\,.\,.\,.\,.\,.\,.} & $0{,}12099$\\[.4em]
\hline
\makebox[17em][l]{$\log\left(\frac{\mathrm{d}b}{\mathrm{d}u}\right)$\hfill .\,.\,.\,.\,.\,.\,.\,.\,.\,.\,.} & $9{,}63241_n$\\[.4em]
\end{tabular}
\end{center}
\begin{center}
\begin{tabular}{@{}ll@{}}
\multicolumn{2}{c}{VII$^{*}$}\\[.4em]
\makebox[17em][l]{$\log r\sin u\cos i$\hfill .\,.} & $9{,}75999_n$\\[.4em]
\makebox[17em][l]{$\log\cos(N-b)$\hfill .\,.} & $9{,}99759_n$\\[.4em]
\makebox[17em][l]{$\mathrm{C.}\log\Delta$\hfill .\,.\,.\,.\,.\,.\,.} & $9{,}91759$\\[.4em]
\makebox[17em][l]{$\mathrm{C.}\log\cos N$\hfill .\,.\,.\,.} & $0{,}00001_n$\\[.4em]
\hline
\makebox[17em][l]{$\log\left(\frac{\mathrm{d}b}{\mathrm{d}i}\right)$\hfill .\,.\,.\,.\,.\,.\,.} & $9{,}67518_n$\\[.4em]
\end{tabular}
\end{center}
\begin{center}
\begin{tabular}{@{}ll@{}}
\multicolumn{2}{c}{VIII}\\[.4em]
\makebox[17em][l]{$({}^{*})$\hfill .\,.\,.\,.\,.\,.\,.\,.\,.\,.\,.} & $9{,}63962$\\[.4em]
\makebox[17em][l]{$\log\sin b\cos b$\hfill .\,.} & $9{,}04212_n$\\[.4em]
\hline
\makebox[17em][l]{$\log\left(\frac{\mathrm{d}b}{\mathrm{d}\AscendingNode}\right)$\hfill .\,.\,.\,.\,.\,.} & $8{,}68174_n$\\[.4em]
\end{tabular}
\end{center}
\par
Collectis hisce valoribus prodit
\[
\begin{aligned}
\mathrm{d}l={}&+0{,}20589\,\mathrm{d}r+1{,}66073\,\mathrm{d}u-0{,}11152\,\mathrm{d}i+1{,}70458\,\mathrm{d}\AscendingNode\\
\mathrm{d}b={}&+0{,}03665\,\mathrm{d}r-0{,}42895\,\mathrm{d}u-0{,}47335\,\mathrm{d}i-0{,}04806\,\mathrm{d}\AscendingNode.
\end{aligned}
\]
Vix necesse erit quod iam saepius monuimus hic repetere, scilicet, vel variationes $\mathrm{d}l$, $\mathrm{d}b$, $\mathrm{d}u$, $\mathrm{d}i$, $\mathrm{d}\AscendingNode$ in partibus radii exprimendas esse, vel coëfficientes ipsius $\mathrm{d}r$ per $206265''$ multiplicandos, si illae in minutis secundis expressae concipiantur.
\par
Designando iam longitudinem perihelii (quae in exemplo nostro est $52^0\,18'\,9''{,}30$) per $\Pi$ atque anomaliam veram per $v$, erit longitudo in orbita $=u+\AscendingNode=v+\Pi$, adeoque $\mathrm{d}u=\mathrm{d}v+\mathrm{d}\Pi-\mathrm{d}\AscendingNode$, quo valore in formulis praecedentibus substituto, $\mathrm{d}l$ et $\mathrm{d}b$ per $\mathrm{d}r$, $\mathrm{d}v$, $\mathrm{d}\Pi$, $\mathrm{d}\AscendingNode$, $\mathrm{d}i$ expressae habebuntur. \SourceLeaf{104}Nihil itaque iam superest, nisi ut $\mathrm{d}r$ et $\mathrm{d}v$ ad normam artt. 15, 16 per variationes differentiales elementorum ellipticorum exhibeantur\footnote{Characterem $M$ in calculo sequente haud amplius angulum nostrum auxiliarem exprimere, sed (ut in Sect. I) anomaliam mediam, quisque sponte videbit.}.
\par
Erat in exemplo nostro, art. 14, $\log\frac{r}{a}=9{,}90355=\log\left(\frac{\mathrm{d}r}{\mathrm{d}a}\right)$
\begin{center}
\begin{tabular}{@{}ll@{}}
\makebox[17em][l]{$\log\frac{aa}{rr}$\hfill .\,.\,.\,.\,.\,.} & $0{,}19290$\\[.4em]
\makebox[17em][l]{$\log\cos\varphi$\hfill .\,.\,.\,.} & $9{,}98652$\\[.4em]
\hline
\makebox[17em][l]{$\log\left(\frac{\mathrm{d}v}{\mathrm{d}M}\right)$\hfill .\,.\,.\,.} & $0{,}17942$\\[.4em]
\end{tabular}
\end{center}
\begin{center}
\begin{tabular}{@{}ll@{}}
$2-e\cos E=$ & $1{,}80085$\\[.4em]
$ee=$ & $0{,}06018$\\[.4em]
\hline
 & $1{,}74067$\\[.4em]
\makebox[17em][l]{$\log$\hfill .\,.\,.\,.\,.\,.\,.\,.} & $0{,}24072$\\[.4em]
\makebox[17em][l]{$\log\frac{aa}{rr}$\hfill .\,.\,.\,.\,.\,.} & $0{,}19290$\\[.4em]
\makebox[17em][l]{$\log\sin E$\hfill .\,.\,.\,.} & $9{,}76634_n$\\[.4em]
\hline
\makebox[17em][l]{$\log\left(\frac{\mathrm{d}v}{\mathrm{d}\varphi}\right)$\hfill .\,.\,.\,.\,.} & $0{,}19996_n$\\[.4em]
\end{tabular}
\end{center}
\begin{center}
\begin{tabular}{@{}ll@{}}
\makebox[17em][l]{$\log a$\hfill .\,.\,.\,.\,.\,.\,.} & $0{,}42244$\\[.4em]
\makebox[17em][l]{$\log\operatorname{tang}\varphi$\hfill .\,.\,.} & $9{,}40320$\\[.4em]
\makebox[17em][l]{$\log\sin v$\hfill .\,.\,.\,.\,.} & $9{,}84931_n$\\[.4em]
\hline
\makebox[17em][l]{$\log\left(\frac{\mathrm{d}r}{\mathrm{d}M}\right)$\hfill .\,.\,.\,.} & $9{,}67495_n$\\[.4em]
\end{tabular}
\end{center}
\begin{center}
\begin{tabular}{@{}ll@{}}
\makebox[17em][l]{$\log a$\hfill .\,.\,.\,.\,.\,.\,.} & $0{,}42244$\\[.4em]
\makebox[17em][l]{$\log\cos\varphi$\hfill .\,.\,.\,.} & $9{,}98652$\\[.4em]
\makebox[17em][l]{$\log\cos v$\hfill .\,.\,.\,.} & $9{,}84966$\\[.4em]
\hline
\makebox[17em][l]{$\log\left(\frac{\mathrm{d}r}{\mathrm{d}\varphi}\right)$\hfill .\,.\,.\,.\,.} & $0{,}25862_n$\\[.4em]
\end{tabular}
\end{center}
\noindent Hinc colligitur
\[
\begin{aligned}
\mathrm{d}v&=+1{,}51154\,\mathrm{d}M-1{,}58475\,\mathrm{d}\varphi\\
\mathrm{d}r&=-0{,}47310\,\mathrm{d}M-1{,}81393\,\mathrm{d}\varphi+0{,}80085\,\mathrm{d}a,
\end{aligned}
\]
quibus valoribus in formulis praecedentibus substitutis prodit
\[
\begin{aligned}
\mathrm{d}l={}&+2{,}41287\,\mathrm{d}M-3{,}00531\,\mathrm{d}\varphi+0{,}16488\,\mathrm{d}a\\
 &+1{,}66073\,\mathrm{d}\Pi-0{,}11152\,\mathrm{d}i+0{,}04385\,\mathrm{d}\AscendingNode\\
\mathrm{d}b={}&-0{,}66572\,\mathrm{d}M+0{,}61331\,\mathrm{d}\varphi+0{,}02935\,\mathrm{d}a\\
 &-0{,}42895\,\mathrm{d}\Pi-0{,}47335\,\mathrm{d}i+0{,}38089\,\mathrm{d}\AscendingNode.
\end{aligned}
\]
\par
Si tempus cui locus computatus respondet $n$ diebus ab epocha distare supponitur, longitudoque media pro epocha per $N$, motus diurnus per $\DailyMotionSign$ denotatur, erit $M=N+n\DailyMotionSign-\Pi$, adeoque $\mathrm{d}M=\mathrm{d}N+n\,\mathrm{d}\DailyMotionSign-\mathrm{d}\Pi$. In exemplo nostro tempus loco computato respondens est Octobris dies $17{,}41501$ anni 1804 sub meridiano Parisiensi: quodsi itaque pro epocha assumitur initium anni 1805, est $n=-74{,}58499$; longitudo media pro epocha ista statuta \SourceLeaf{105}fuerat $=41^0\,52'\,21''{,}68$, motusque diurnus $=824''{,}7990$. Substituto iam in formulis modo inventis pro $\mathrm{d}M$ valore suo, mutationes differentiales loci geocentrici per solas mutationes elementorum expressae ita se habent:
\[
\begin{aligned}
\mathrm{d}l={}&+2{,}41287\,\mathrm{d}N-179{,}96\,\mathrm{d}\DailyMotionSign-0{,}75214\,\mathrm{d}\Pi\\
 &-3{,}00531\,\mathrm{d}\varphi+0{,}16488\,\mathrm{d}a-0{,}11152\,\mathrm{d}i+0{,}04385\,\mathrm{d}\AscendingNode\\
\mathrm{d}b={}&-0{,}66572\,\mathrm{d}N+49{,}65\,\mathrm{d}\DailyMotionSign+0{,}23677\,\mathrm{d}\Pi\\
 &+0{,}61331\,\mathrm{d}\varphi+0{,}02935\,\mathrm{d}a-0{,}47335\,\mathrm{d}i+0{,}38089\,\mathrm{d}\AscendingNode.
\end{aligned}
\]
\par
Si corporis coelestis massa vel negligitur vel saltem tamquam cognita spectatur, $\DailyMotionSign$ et $a$ ab invicem dependentes erunt, adeoque vel $\mathrm{d}\DailyMotionSign$ vel $\mathrm{d}a$ e formulis nostris eliminare licebit. Scilicet quum per art. 6 habeatur $\DailyMotionSign a^{\frac{3}{2}}=k\BareRoot{(1+\mu)}$, erit $\frac{\mathrm{d}\DailyMotionSign}{\DailyMotionSign}=-\frac{3}{2}\frac{\mathrm{d}a}{a}$, in qua formula, si $\mathrm{d}\DailyMotionSign$ in partibus radii exprimenda est, etiam $\DailyMotionSign$ perinde exprimere oportebit. Ita in exemplo nostro habetur
\begin{center}
\begin{tabular}{@{}ll@{}}
\makebox[17em][l]{$\log\DailyMotionSign$\hfill .\,.\,.\,.\,.} & $2{,}91635$\\[.4em]
\makebox[17em][l]{$\log1''$\hfill .\,.\,.\,.} & $4{,}68557$\\[.4em]
\makebox[17em][l]{$\log\frac{3}{2}$\hfill .\,.\,.\,.\,.} & $0{,}17609$\\[.4em]
\makebox[17em][l]{$\mathrm{C.}\log a$\hfill .\,.\,.} & $9{,}57756$\\[.4em]
\hline
\makebox[17em][l]{$\log\frac{\mathrm{d}\DailyMotionSign}{\mathrm{d}a}$\hfill .\,.\,.\,.} & $7{,}35557_n$\\[.4em]
\end{tabular}
\end{center}
\noindent sive $\mathrm{d}\DailyMotionSign=-0{,}0022676\,\mathrm{d}a$, atque $\mathrm{d}a=-440{,}99\,\mathrm{d}\DailyMotionSign$, quo valore in formulis nostris substituto, tandem emergit forma ultima:
\[
\begin{aligned}
\mathrm{d}l={}&+2{,}41287\,\mathrm{d}N-252{,}67\,\mathrm{d}\DailyMotionSign-0{,}75214\,\mathrm{d}\Pi\\
 &-3{,}00531\,\mathrm{d}\varphi-0{,}11152\,\mathrm{d}i+0{,}04385\,\mathrm{d}\AscendingNode\\
\mathrm{d}b={}&-0{,}66572\,\mathrm{d}N+36{,}71\,\mathrm{d}\DailyMotionSign+0{,}23677\,\mathrm{d}\Pi\\
 &+0{,}61331\,\mathrm{d}\varphi-0{,}47335\,\mathrm{d}i+0{,}38089\,\mathrm{d}\AscendingNode.
\end{aligned}
\]
In evolutione harum formularum omnes mutationes $\mathrm{d}l$, $\mathrm{d}b$, $\mathrm{d}N$, $\mathrm{d}\DailyMotionSign$, $\mathrm{d}\Pi$, $\mathrm{d}\varphi$, $\mathrm{d}i$, $\mathrm{d}\AscendingNode$ in partibus radii expressas supposuimus, manifesto autem propter homogeneitatem omnium partium eaedem formulae etiamnum valebunt, si omnes illae mutationes in minutis secundis exprimuntur.
\begin{center}\rule{7em}{.4pt}\end{center}
